\documentclass[12pt]{article}
\usepackage{amsmath}
\usepackage{graphicx}
\usepackage{algorithm}
\usepackage{algpseudocode}
\usepackage{caption}

\title{A Study on Computational Metrics}
\author{Bodike Chaithali \\ B24CS013}
\date{\today}

\begin{document}
	\maketitle
	
	\begin{abstract}
		This report demonstrates essential \LaTeX{} features including structured
		sections, equations, tables, figures, algorithms, custom macros, and bibliography
		management. We illustrate their usage through a new context with different formulas.
	\end{abstract}
	
	\section{Introduction}
	Mathematical relationships are fundamental in computational studies. 
	For example, the Gaussian integral is given by
	\begin{equation}
		\int_{-\infty}^{\infty} e^{-x^2}\,dx = \sqrt{\pi},
	\end{equation}
	which we will reference in later discussions.
	
	\section{Methods}
	We measured the refractive index of the material using the custom macro
	\begin{equation}
		n = \frac{c}{v},
	\end{equation}
	where $c$ is the speed of light in vacuum and $v$ is the speed of light in the material.
	
	\subsection{Figures}
	Figure~\ref{fig:sample} shows a demonstration image included to illustrate figure handling.
	
	\begin{figure}[h]
		\centering
		\includegraphics[width=0.4\textwidth]{"/home/bodike/Desktop/Sem3/CSP203/Practice/LaTeX_logo.png"} 
		\caption{Sample illustration figure.}
		\label{fig:sample}
	\end{figure}
	
	\subsection{Tables}
	Table~\ref{tab:readings} shows sample measurements of sensor readings.
	
	\begin{table}[h]
		\centering
		\begin{tabular}{|c|c|c|c|}
			\hline
			Sensor & Reading 1 & Reading 2 & Reading 3 \\
			\hline
			A & 12.5 & 13.1 & 12.9 \\
			\hline
			B & 14.0 & 13.8 & 14.2 \\
			\hline
			C & 11.8 & 12.0 & 11.9 \\
			\hline
		\end{tabular}
		\caption{Sample sensor readings.}
		\label{tab:readings}
	\end{table}
	\section{Algorithm}
	Algorithm~\ref{alg:factorial} demonstrates computing the factorial of a number.
	
	\begin{algorithmic}
		\State $n \gets$ input integer
		\State $result \gets 1$
		\For{$i \gets 1$ \textbf{to} $n$}
		\State $result \gets result \times i$
		\EndFor
	\end{algorithmic}
	\captionof{algorithm}{Compute factorial of an integer.}
	\label{alg:factorial}
\end{document}